\documentclass[12pt]{amsart}

%% ─── Packages ────────────────────────────────────────────────────────────────
\usepackage{amsmath,amssymb,amsthm}
\usepackage{geometry}
\geometry{margin=1in}
\usepackage{hyperref}
\hypersetup{
  colorlinks=true,
  linkcolor=blue,
  citecolor=blue,
  urlcolor=blue
}
\usepackage{booktabs}
\usepackage{microtype}
\usepackage{enumitem}

%% ─── Theorem Environments ────────────────────────────────────────────────────
\theoremstyle{plain}
\newtheorem{theorem}{Theorem}[section]
\newtheorem{proposition}[theorem]{Proposition}
\newtheorem{lemma}[theorem]{Lemma}

\theoremstyle{definition}
\newtheorem{definition}[theorem]{Definition}
\newtheorem{example}[theorem]{Example}
\newtheorem{criterion}[theorem]{Criterion}

\theoremstyle{remark}
\newtheorem{remark}[theorem]{Remark}
\newtheorem{principle}[theorem]{Principle}

%% ─── Metadata ────────────────────────────────────────────────────────────────
\title[Spherical Cow Philosophy]{Spherical Cow Philosophy:
  A Methodological Framework for Mathematical Discovery}

\author{Khayyam Wakil}
\address{The ARC Institute of Knowware, Calgary, Alberta, Canada}
\email{the@knowware.institute}
\date{February 2026}

\keywords{mathematical methodology, constitutional forcing, spherical cow,
  empirical constitution, twin primes, sieve theory, methodological pluralism,
  structural barriers, cross-domain validation}

\subjclass[2020]{
  00A30,  % Philosophy of mathematics
  00A35,  % Methodology of mathematics
  11N35,  % Sieves
  11N36   % Applications of sieve methods
}

%% ─── Abstract ────────────────────────────────────────────────────────────────
\begin{document}

\begin{abstract}
We present \emph{spherical cow philosophy}---a methodological framework for
mathematical discovery that complements classical proof-by-contradiction by
accepting empirical constitutional constraints as foundational and deriving their
forced consequences. Named after the physics heuristic of modelling complex
systems through essential simplification, the approach inverts the usual
hierarchy: instead of deriving special cases from universal axioms, we identify
structural constraints that are empirically universal within a domain, derive
what they force rigorously, and validate predictions across independent domains.

The framework addresses a specific situation: when classical sieve methods
encounter persistent structural barriers---such as the $\theta = 1/2$ ceiling
for twin primes over two centuries of effort---the productive path may be to
accept the constitutional constraint (here, $p \equiv 2 \pmod{3}$ for all twin
primes $p > 3$) as the natural entry point and derive its arithmetic consequences,
rather than to approach the barrier as an obstacle for more powerful general
methods.

Applied to twin primes, this methodology motivates the cascade moduli framework
proved in the companion paper~\cite{WakilEisenstein}, achieves level of
distribution $\theta_W = 5/8$ for cascade moduli $q = 3^K$, and connects the
arithmetic result to a structural parallel in information theory documented
in~\cite{WakilShannon}. The combinatorial foundation is traced to the mod-$3$
fractal of Khayyam's Triangle~\cite{WakilKhayyam}. The framework is explicitly
falsifiable: it predicts a confirmed non-instance (Sophie Germain primes at
$\theta = 1/2$) and an open prediction ($k = 4$ constellation at $\theta = 3/4$)
via the $k$-hierarchy of~\cite{WakilShannon}.

We establish precise criteria for when the framework applies, contrast it with
classical approaches, and distinguish it carefully from circular or post-hoc
reasoning. Logical status labels are attached to every claim.
\end{abstract}

\maketitle
\tableofcontents

%% ─────────────────────────────────────────────────────────────────────────────
\section{Introduction: The Physics Analogy}
%% ─────────────────────────────────────────────────────────────────────────────

\subsection{The original spherical cow}

When physicists model heat transfer from a cow, they make a radical
simplification: assume a spherical cow. Strip away legs, spots, tail, digestive
complexity---keep only mass distribution and thermal properties. As
Clark~\cite{Clark2020} articulates:
\begin{quote}
Make it as simple as possible in order to take your understanding beyond the
simple\ldots Without saying ``a cow is a sphere''---an imperfect, simple
description---we would never have been in the position to describe it accurately
at all.
\end{quote}
The spherical cow succeeds because it knows what to \emph{ignore}. Domain-specific
details obscure the essential physics. By ignoring the irrelevant, physicists
isolate what matters, solve the simplified problem, and add complexity back
strategically where it is needed.

\subsection{The mathematical parallel}

Classical mathematics has three primary responses to a hard problem:
\begin{enumerate}[label=(\arabic*)]
  \item Assume the statement false, derive consequences, seek a contradiction.
  \item Build from axioms: prove the statement from more fundamental principles.
  \item Exhaust known techniques until something works.
\end{enumerate}
These work brilliantly when the problem is provable from accepted axioms.
When they encounter a persistent structural barrier, the barrier itself is
information. It may be signalling that the productive formulation has not yet
been identified---not that the problem is harder than expected.

\emph{Spherical cow philosophy} provides a complementary response: instead of
asking ``how do we overcome this barrier?'', ask ``what constraint is this
barrier telling us is constitutional?'' Accept that constraint as the productive
entry point. Derive what it forces. Validate the consequences.

\subsection{Position in the program}

This paper is the methodological companion to a five-paper arithmetic
program~\cite{WakilKhayyam,WakilRamanujan,WakilEisenstein,WakilW3,WakilShannon}.

\begin{itemize}
  \item \cite{WakilKhayyam} establishes the historical and combinatorial
    foundation: Khayyam's Triangle mod~$3$ and the fractal structure that
    underlies the cascade moduli framework.
  \item \cite{WakilRamanujan} derives the universal formula
    $\theta_k = (2^k - k)/2^k$ from combinatorial first principles and proves
    it analytically for $k = 3$.
  \item \cite{WakilEisenstein} establishes $\theta_W = 5/8$ for cascade
    moduli $q = 3^K$ via the geometry of the Eisenstein integers $\mathbb{Z}[\omega]$.
  \item \cite{WakilW3} verifies the Siegel zero condition (Condition~W3)
    required by~\cite{WakilEisenstein} via the CM structure of $\mathbb{Q}(\omega)$.
  \item \cite{WakilShannon} documents the structural parallel between the
    arithmetic result and Shannon's theorems in information theory, and states
    the falsifiable $k$-hierarchy predictions.
\end{itemize}

The present paper explains \emph{why} that program is structured the way it is,
and what methodological principle makes it natural and non-arbitrary.

\subsection{What this paper claims and does not claim}
\label{sec:claims}

We attach logical-status labels to every significant claim throughout.
Here we state the top-level ones explicitly.

\begin{itemize}
  \item[\textbf{[Claimed]}] Accepting $p \equiv 2 \pmod{3}$ as the constitutional
    entry point for twin primes is methodologically productive: it motivates the
    cascade moduli framework and leads to the proved result $\theta_W = 5/8$
    of~\cite{WakilEisenstein}.

  \item[\textbf{[Claimed]}] The constitution $p \equiv 2 \pmod{3}$ is a theorem
    (provable in two lines), not an axiom. Accepting a proved fact as the
    \emph{productive entry point}---rather than a consequence to be derived
    from something deeper---is the methodological move this paper advocates.

  \item[\textbf{[Claimed]}] The framework is falsifiable. The Sophie Germain
    non-instance (Section~\ref{sec:noncircular}) was predicted and confirmed.
    The $k=4$ prediction (Section~\ref{sec:decisioncriteria}) is open and testable.

  \item[\textbf{[Not claimed]}] That constitutional forcing alone proves twin
    prime infinitude. The result $\theta_W = 5/8$ is established for cascade
    moduli $q = 3^K$; extending this to the full moduli input required by
    GPY~\cite{GPY2009} is the principal open step identified
    in~\cite{WakilEisenstein}.

  \item[\textbf{[Not claimed]}] That twin prime infinitude is independent of
    ZFC in the technical G\"odel--Cohen sense. We make no claim about
    independence. The methodology applies whenever constitutional structure is
    the productive entry point, regardless of the problem's axiomatic status.
\end{itemize}

%% ─────────────────────────────────────────────────────────────────────────────
\section{The Classical Toolkit and Its Limits}
%% ─────────────────────────────────────────────────────────────────────────────

\subsection{When classical methods succeed}

Proof by contradiction is among the most powerful tools in mathematics.

\begin{example}[Irrationality of $\sqrt{2}$]
Assume $\sqrt{2} = p/q$ in lowest terms. Then $2q^2 = p^2$, forcing $p$ even;
writing $p = 2k$ gives $q^2 = 2k^2$, forcing $q$ even---contradicting lowest
terms. Therefore $\sqrt{2}$ is irrational.
\end{example}

This works because the statement is provable from the axioms of arithmetic:
the assumption of rationality produces a genuine logical contradiction.

\begin{example}[Infinitude of primes, Euclid]
Suppose there are finitely many primes $p_1, \ldots, p_n$. The number
$N = p_1 \cdots p_n + 1$ is divisible by none of them, so it has a prime factor
not in the list---a contradiction.
\end{example}

Both examples share a structure: \emph{the assumption of falsehood leads
directly to a contradiction derivable from accepted axioms}. Classical proof
by contradiction works exactly when this structure is present.

\subsection{Where structural barriers arise}
\label{sec:barriers}

Some problems resist sustained classical effort not because they are
unprovable, but because the productive formulation has not yet been identified.
The barrier is information: it is telling us something about the structure of
the problem.

\begin{example}[The $\theta = 1/2$ barrier for twin primes]
\label{ex:twinprime-barrier}
The Goldston--Pintz--Y{\i}ld{\i}r{\i}m theorem~\cite{GPY2009} shows that a
level of distribution $\theta > 1/2$ for primes in arithmetic progressions
would imply the existence of infinitely many prime pairs with bounded gaps,
and in particular---via a refinement---would yield infinitely many twin primes.
Classical Bombieri--Vinogradov~\cite{BV1965} gives $\theta = 1/2$ for
\emph{all} moduli. Zhang~\cite{Zhang2014} achieved $\theta = 1/2 + \delta$
for a small $\delta > 0$ and smooth moduli, proving bounded gaps. Maynard and
Tao~\cite{Maynard2015} improved the gap bound substantially. But the gap
$= 2$ (twin primes specifically) remained out of reach.

The barrier at $\theta = 1/2$ for general moduli is not arbitrary. It reflects
the limits of the large sieve in full generality. Classical methods treat all
moduli uniformly; they cannot exploit structure specific to a particular prime.
\end{example}

The spherical cow response to Example~\ref{ex:twinprime-barrier} is not to
work harder on the general barrier. It is to ask: \emph{what constitutional
structure do twin primes carry that distinguishes them from the general case?}

The answer---$p \equiv 2 \pmod{3}$ for all twin primes $p > 3$---is a two-line
theorem. Accepting it as the productive entry point leads immediately to cascade
moduli, Eisenstein integer geometry, and $\theta_W = 5/8$
(see~\cite{WakilRamanujan,WakilEisenstein}).

\subsection{The structural limitation of the classical approach}
\label{sec:limitation}

Classical methods seek \emph{universal} patterns: results that hold for all
moduli, all primes, all arithmetic progressions. This universality is their
strength and their limitation. When a problem has constitutional structure that
singles out a specific prime (here, $3$), universal methods must either ignore
that structure or work much harder to recover it.

Spherical cow philosophy embraces the specific structure from the outset. It
trades universality for depth, and gains access to geometric tools---the
hexagonal lattice of $\mathbb{Z}[\omega]$, the CM field $\mathbb{Q}(\omega)$,
Kloosterman cancellation in a cyclic group---that are invisible from the
universal vantage point.

%% ─────────────────────────────────────────────────────────────────────────────
\section{Spherical Cow Philosophy: The Framework}
%% ─────────────────────────────────────────────────────────────────────────────

\subsection{The inversion}

Classical hierarchy:
\[
\text{Universal Axioms}
\;\longrightarrow\;
\text{General Theorems}
\;\longrightarrow\;
\text{Special Cases}
\;\longrightarrow\;
\text{[barrier]}
\]
Spherical cow hierarchy:
\begin{align*}
&\text{Empirical Constitution}
\;\longrightarrow\;
\text{Forced Consequences}\\
&\quad\longrightarrow\;
\text{Cross-Domain Validation}
\;\longrightarrow\;
\text{Falsifiable Predictions}
\end{align*}

The key difference: we do not try to prove the constitution from something
``more basic.'' We show the constitution generates observed structure across
multiple independent domains and makes predictions that can be confirmed or
refuted.

\subsection{The five steps}
\label{sec:fivesteps}

\begin{definition}[Spherical Cow Methodology]
\label{def:methodology}
Given a domain with a persistent structural barrier, the spherical cow
methodology proceeds as follows.

\smallskip
\noindent\textbf{Step 1: Identify the empirical constitution.}
Find a constraint that is:
\begin{enumerate}[label=(\alph*)]
  \item \emph{Verifiable by direct computation}---not an assumption;
  \item \emph{Universal within the domain}---all instances satisfy it;
  \item \emph{Simple to state}---not complex conditional logic.
\end{enumerate}

\noindent\textbf{Step 2: Derive forced consequences.}
Using rigorous mathematics, determine what the constitutional constraint
\emph{forces}. Do not add hypotheses; derive consequences.

\noindent\textbf{Step 3: Validate within the domain.}
Check that the forced consequences match domain observations, both
computationally and analytically where possible.

\noindent\textbf{Step 4: Validate across independent domains.}
Determine whether the same constitutional structure explains phenomena in
genuinely independent fields. Independent confirmation strengthens the case
that the constitution is fundamental, not accidental.

\noindent\textbf{Step 5: Generate falsifiable predictions.}
Identify consequences of the constitution that have not yet been observed,
and check whether they hold or fail. A framework that makes and confirms
predictions---and identifies genuine non-instances---is more than decorative.
\end{definition}

\subsection{The spherical cow principle}
\label{sec:principle}

\begin{principle}[Spherical Cow]
\label{prin:sc}
What to ignore: domain-specific complexity that obscures constitutional
structure.\\
What to keep: the constitutional constraint that generates forced consequences.\\
How to validate: cross-domain prediction and confirmed non-instances.
\end{principle}

The physics analogy is precise. A physicist ignores the cow's legs and spots
to isolate thermal properties. A mathematician ignores the full generality of
all moduli to isolate the constitutional property of twin primes ($p \equiv 2
\pmod{3}$) and the cascade structure it forces.

%% ─────────────────────────────────────────────────────────────────────────────
\section{Case Study: Twin Primes}
%% ─────────────────────────────────────────────────────────────────────────────

\subsection{The constitution}
\label{sec:constitution}

\begin{proposition}[Twin Prime Constitution]
\label{prop:constitution}
Every twin prime $p > 3$ satisfies $p \equiv 2 \pmod{3}$.
\end{proposition}

\begin{proof}
Among any three consecutive integers, exactly one is divisible by~$3$.
For the pair $(p, p+2)$: if $p \equiv 0 \pmod{3}$ then $3 \mid p$, so $p$
is not prime (for $p > 3$). If $p \equiv 1 \pmod{3}$ then $p + 2 \equiv 0
\pmod{3}$, so $3 \mid p+2$, hence $p+2$ is not prime (for $p+2 > 3$). Therefore
$p \equiv 2 \pmod{3}$.
\end{proof}

\begin{remark}
Proposition~\ref{prop:constitution} is a theorem, not an axiom. The
methodological move is not to assume it but to \emph{accept it as the productive
entry point}---to ask what arithmetic structure it forces, rather than treating
it as a consequence to be recovered from something more general.
\end{remark}

\subsection{The forced consequences}
\label{sec:forced}

Given Proposition~\ref{prop:constitution}, the spherical cow methodology asks:
what does $p \equiv 2 \pmod{3}$ force?

\begin{enumerate}[label=(\arabic*)]
  \item \textbf{Cascade structure.} By the Chinese Remainder Theorem,
    $\gcd(3^K, \ell) = 1$ for all primes $\ell \neq 3$, so modular conditions
    modulo $3^K$ and modulo $\ell$ are independent. The constitutional condition
    forces attention to the specific family of moduli $q = 3^K$---the
    \emph{cascade moduli}.

  \item \textbf{Eisenstein geometry.} The prime $3$ ramifies in the Eisenstein
    integers $\mathbb{Z}[\omega]$ as $(3) = (1-\omega)^2$ (Proposition~2.2
    of~\cite{WakilEisenstein}). A twin prime $p \equiv 2 \pmod{3}$ is
    \emph{inert} in $\mathbb{Z}[\omega]$; its companion $p+2 \equiv 1 \pmod{3}$
    \emph{splits}. Twin primes are precisely the inert-split pairs of the
    hexagonal lattice (Corollary~6.4 of~\cite{WakilEisenstein}).

  \item \textbf{Three-level filtration.} The ramification $(3) = (1-\omega)^2$
    induces a natural filtration on $\mathbb{Z}[\omega]/(3^K)$ with exactly
    three independent levels, forcing sieve weights over cascade moduli to be
    triply-well-factorable (Corollary~2.5 of~\cite{WakilEisenstein}).

  \item \textbf{Level of distribution $\theta_W = 5/8$.} Three-fold
    Cauchy--Schwarz over the filtration gives ceiling $\theta \leq 2/3$;
    cyclic group cancellation from $({\mathbb{Z}}/3^K{\mathbb{Z}})^* \cong
    \mathbb{Z}/(2 \cdot 3^{K-1}\mathbb{Z})$ improves this to $\theta_W = 5/8$
    (Theorem~1.1 of~\cite{WakilEisenstein}, with Condition~W3 verified
    in~\cite{WakilW3}).

  \item \textbf{Combinatorial origin.} The value $5/8 = (2^3 - 3)/2^3$
    counts valid configurations in a three-level constitutional sieve: of the
    $2^3 = 8$ binary configurations, exactly $k = 3$ are invalid, leaving $5$
    valid (Theorem~3.4 of~\cite{WakilRamanujan}). This count traces to the
    mod-$3$ Sierpi\'nski structure of Khayyam's Triangle
    (Proposition~2.3 of~\cite{WakilRamanujan}, and~\cite{WakilKhayyam}).
\end{enumerate}

All five consequences are forced by the single constitutional fact of
Proposition~\ref{prop:constitution}. None requires additional hypotheses.

\subsection{The GPY connection and the open step}
\label{sec:gpy}

The Goldston--Pintz--Y{\i}ld{\i}r{\i}m theorem~\cite{GPY2009} requires a level
of distribution $\theta > 1/2$ over a sufficiently broad family of moduli to
conclude twin prime infinitude. The result $\theta_W = 5/8$ established
in~\cite{WakilEisenstein} applies to cascade moduli $q = 3^K$: a one-parameter
family indexed by powers of a single prime.

\textbf{[Open]} Extending $\theta_W = 5/8$ to a moduli family broad
enough to serve as GPY input is the principal open step in this program.
Two avenues are identified in~\cite{WakilEisenstein}: (a) combining cascade
moduli with smooth moduli via Pascadi's result~\cite{Pascadi2025}; (b) developing
a GPY variant that works with the specific cascade family directly.

We do not claim this step is complete. The spherical cow methodology motivates
and structures the arithmetic program; the completion of the program is the work
of the companion papers and their successors.

\subsection{Cross-domain validation: Shannon's theorems}
\label{sec:shannon}

The same constitutional structure---an exponential space forced by an intrinsic
constraint onto a strict subspace with a uniquely determined governing
constant---appears in Shannon's two foundational theorems of 1948~\cite{Shannon1948},
as documented in detail in~\cite{WakilShannon}.

\begin{itemize}
  \item \textbf{Asymptotic Equipartition Property.} $|\mathcal{A}|^n$ sequences
    concentrate onto a typical set of size $\approx 2^{nH}$, where entropy $H$
    is uniquely forced by the source distribution.
  \item \textbf{Channel Coding Theorem.} Reliable communication requires rate
    $R < C$, where capacity $C$ is uniquely forced by the channel.
\end{itemize}

Both are instances of the same pattern as the cascade moduli result: a structural
constraint forces an exponential space to concentrate onto a subspace, governed
by a constant that is uniquely determined by the constraint and cannot be altered
without changing the underlying structure. The formal definition and proof that
all three instances satisfy the same conditions are given in~\cite{WakilShannon}.

The 77-year gap between Shannon's results (1948) and the arithmetic instance
(2026) is not surprising: the connection is visible only after both sides are
independently established. This order---independent results first, structural
comparison second---guarantees that neither result was constructed to match the
other.

%% ─────────────────────────────────────────────────────────────────────────────
\section{Non-Circularity and Falsifiability}
%% ─────────────────────────────────────────────────────────────────────────────

\subsection{Why this is not circular}
\label{sec:noncircular}

A potential objection: the framework accepts the constitution, derives
consequences, and validates against observations in the same domain. Is this
circular?

It is not, for three reasons.

\textbf{First}, the constitution is a theorem (Proposition~\ref{prop:constitution}),
not an assumption. Accepting a proved fact as the entry point is not circular;
it is a choice of direction.

\textbf{Second}, the consequences are derived by rigorous mathematics (CRT,
Eisenstein geometry, Cauchy--Schwarz, Watt's theorem), not assumed. The
derivation could in principle have produced a different value of $\theta$, or
no useful bound at all. The fact that it produces $5/8 > 1/2$ is a non-trivial
output.

\textbf{Third}, and most importantly, the framework makes \emph{falsifiable
predictions that were confirmed before the verification was run}. The Sophie
Germain non-instance is the clearest example.

\begin{proposition}[Sophie Germain Non-Instance, {\cite[Theorem~7.1]{WakilShannon}}]
\label{prop:sg}
Sophie Germain prime pairs $(p, 2p+1)$ have dimensional parameter $k = 2$: the
mod-$3$ residue of $2p+1$ is determined by that of $p$ (both satisfy
$p \equiv 2 \pmod{3}$), giving only two independent constraint levels rather
than three. The framework therefore predicts level of distribution
$\theta_2 = (2^2 - 2)/2^2 = 1/2$, with no improvement over Bombieri--Vinogradov.
Computational analysis confirms this: Sophie Germain primes do not exhibit
the cascade moduli enhancement.
\end{proposition}

\begin{remark}
The Sophie Germain prediction follows the same logic as the twin prime result:
coupled constraints reduce the effective dimensionality. In information theory,
a degraded channel has lower capacity because input-output dependence increases;
in Sophie Germain primes, the coupling between $p$ and $2p+1$ modulo~$3$
reduces effective $k$ from $3$ to $2$. The framework handles both cases with
the same formula $\theta_k = (2^k - k)/2^k$, and correctly predicts which
constellations do and do not exhibit the enhancement.
\end{remark}

\subsection{Decision criteria: when to use this framework}
\label{sec:decisioncriteria}

\begin{criterion}[Apply spherical cow philosophy when:]
\label{crit:apply}
\begin{enumerate}[label=(\arabic*)]
  \item \textbf{Structural barrier}: Classical methods have encountered a
    persistent ceiling that is not merely a technical obstacle but appears to
    reflect the limits of the general approach.
  \item \textbf{Identifiable constitution}: A constraint exists that is
    universal within the domain, verifiable computationally, and simple to
    state.
  \item \textbf{Cross-domain echo}: The same structural pattern appears in
    independent fields, suggesting the constitution captures something
    fundamental rather than domain-specific.
  \item \textbf{Falsifiable predictions}: The framework generates predictions
    about cases not yet examined, which can confirm or refute it.
\end{enumerate}
\end{criterion}

\begin{criterion}[Prefer classical methods when:]
\label{crit:classical}
\begin{enumerate}[label=(\arabic*)]
  \item The statement is provable from accepted axioms by known techniques.
  \item The problem is single-domain without structural echoes elsewhere.
  \item No universal constitutional constraint is identifiable.
  \item Proof by contradiction successfully finds contradictions.
\end{enumerate}
\end{criterion}

\begin{remark}[The $k=4$ open prediction]
\label{rem:k4}
The most immediate falsifiable prediction of this framework is the existence of
a prime constellation with $k = 4$ independent modular constraints achieving
level of distribution $\theta_4 = (2^4 - 4)/2^4 = 3/4$
(see Table~3 of~\cite{WakilShannon}). Identifying such a constellation and
verifying or disproving $\theta_4 = 3/4$ computationally is a concrete open
problem that would either strengthen or refute the $k$-hierarchy.
\end{remark}

%% ─────────────────────────────────────────────────────────────────────────────
\section{Philosophical Foundations}
%% ─────────────────────────────────────────────────────────────────────────────

\subsection{Empiricism and rationalism in mathematics}

Classical mathematics is rationalist: truth derives from axioms and logical
deduction. Empirical observations serve as heuristics, not proofs.

Spherical cow philosophy is \emph{methodologically} empiricist, without
abandoning rigour. Constitutional constraints are accepted as entry points
because they are empirically universal and computationally verifiable. The
\emph{derivation} from those constraints remains fully rigorous. The shift is
in what counts as the productive starting point, not in what counts as a proof.

The physics analogy is again precise. Physicists accept the speed of light $c$
as an empirical constant and rigorously derive the structure of spacetime from
it. The derivation (special relativity, $E = mc^2$, time dilation) is
mathematically exact; the starting point is empirical. Spherical cow philosophy
proposes the same division of labour for mathematics: accept the constitutional
constraint empirically, derive its consequences rigorously.

\subsection{G\"odel, independence, and the scope of this claim}
\label{sec:godel}

G\"odel's incompleteness theorems~\cite{Godel1931} establish that any consistent
formal system containing arithmetic has true statements unprovable within the
system. Cohen~\cite{Cohen1963} demonstrated that the continuum hypothesis is
independent of ZFC.

We do \emph{not} claim that twin prime infinitude is independent of ZFC.
We make no claim about the axiomatic status of any conjecture in this program.

What we do claim is narrower and more useful: \emph{for problems where classical
methods have encountered structural barriers, accepting an empirical constitution
as the entry point is a productive methodological choice}. This claim is
independent of whether the problem is independent of ZFC. It applies whenever
the constitutional structure is the productive formulation, regardless of the
problem's status in formal logic.

\subsection{Methodological pluralism}

Mathematics benefits from multiple methodologies:
\begin{itemize}
  \item Classical proof by contradiction: for statements provable from axioms.
  \item Constructive mathematics: for explicit object construction.
  \item Probabilistic methods: for deterministic results via probability
    (Erd\H{o}s).
  \item Computational verification: for finite cases (four-colour theorem).
  \item \emph{Spherical cow philosophy}: for structural barriers with
    identifiable empirical constitutions and cross-domain echoes.
\end{itemize}

These are not competing methodologies. They are tools for different situations.
The contribution of this paper is to articulate the conditions under which the
last tool is the right one, and to demonstrate its application through the
arithmetic program of~\cite{WakilRamanujan,WakilEisenstein,WakilW3,WakilShannon}.

%% ─────────────────────────────────────────────────────────────────────────────
\section{Conclusion}
%% ─────────────────────────────────────────────────────────────────────────────

\subsection{Summary}

Spherical cow philosophy provides a methodological complement to classical
proof-by-contradiction for problems where structural barriers persist. By
accepting empirical constitutional constraints as the productive entry
point---deriving forced consequences rigorously and validating predictions
across independent domains---the framework opens paths that the classical
universal approach cannot see.

The key results demonstrated in this program:
\begin{enumerate}[label=(\arabic*)]
  \item The constitution $p \equiv 2 \pmod{3}$ forces cascade structure, Eisenstein
    geometry, and level of distribution $\theta_W = 5/8$ for cascade
    moduli~\cite{WakilEisenstein,WakilW3}.
  \item The same structural pattern---exponential space forced onto a strict
    subspace by an intrinsic constraint---appears in Shannon's AEP and Channel
    Coding Theorem, documented in~\cite{WakilShannon}.
  \item The Sophie Germain non-instance confirms the framework's predictive
    power: coupled constraints reduce effective dimensionality exactly as the
    formula predicts.
  \item The $k=4$ open prediction provides a concrete falsifiability test.
\end{enumerate}

\subsection{The philosophical shift}

From: mathematics begins with axioms and proceeds by deduction.

To: mathematics also begins with constitutional constraints and proceeds by
rigorous derivation of their forced consequences.

This is not a rejection of rigour. The derivations in~\cite{WakilEisenstein,
WakilW3} are as rigorous as any analytic number theory. The shift is in what
counts as the \emph{starting point}, not in what counts as a \emph{proof}.

\subsection{Final reflection}

Omar Khayyam spent his life finding geometry hidden in numbers and in the sky.
His calendar was accurate to within a second per year, computed with naked-eye
instruments and geometric sight~\cite{WakilKhayyam}. His triangle mod~$3$
contains the fractal that underlies the cascade moduli framework. He did not
overcome the complexity of the solar year or of cubic equations by working harder
within established methods. He found a different entry point.

The spherical cow reveals what matters. Khayyam knew this. The companion papers
make it rigorous.

%% ─────────────────────────────────────────────────────────────────────────────
\section*{Acknowledgements}
%% ─────────────────────────────────────────────────────────────────────────────

The author thanks Claude Shannon posthumously: his 1948 paper established two
instances of the Shannon--Wakil effect before either the effect or the arithmetic
instance were identified. The author thanks Simon Clark for articulating the
spherical cow principle with the clarity that made the mathematical parallel
visible.

%% ─────────────────────────────────────────────────────────────────────────────
\begin{thebibliography}{99}
%% ─────────────────────────────────────────────────────────────────────────────

%% ── Foundational references ───────────────────────────────────────────────────

\bibitem{Clark2020}
S.~Clark,
\emph{Assume a Spherical Cow: A Physicist's Guide to Understanding the
Universe},
Quadrille Publishing, 2020.

\bibitem{Shannon1948}
C.~E.~Shannon,
A mathematical theory of communication,
\emph{Bell System Tech.\ J.}\ \textbf{27} (1948), 379--423, 623--656.

\bibitem{Godel1931}
K.~G\"odel,
\"Uber formal unentscheidbare S\"atze der Principia Mathematica und verwandter
Systeme~I,
\emph{Monatsh.\ Math.\ Phys.}\ \textbf{38} (1931), 173--198.

\bibitem{Cohen1963}
P.~J.~Cohen,
The independence of the continuum hypothesis,
\emph{Proc.\ Nat.\ Acad.\ Sci.\ USA}\ \textbf{50} (1963), 1143--1148.

%% ── Sieve theory ──────────────────────────────────────────────────────────────

\bibitem{BV1965}
E.~Bombieri and A.~I.~Vinogradov,
On the large sieve,
\emph{Mathematika}\ \textbf{12} (1965), 201--225.

\bibitem{GPY2009}
D.~A.~Goldston, J.~Pintz, and C.~Y.~Y{\i}ld{\i}r{\i}m,
Primes in tuples~I,
\emph{Ann.\ of Math.}\ \textbf{170} (2009), 819--862.

\bibitem{Zhang2014}
Y.~Zhang,
Bounded gaps between primes,
\emph{Ann.\ of Math.}\ \textbf{179} (2014), 1121--1174.

\bibitem{Maynard2015}
J.~Maynard,
Small gaps between primes,
\emph{Ann.\ of Math.}\ \textbf{181} (2015), 383--413.

\bibitem{Pascadi2025}
A.~Pascadi,
Level of distribution $5/8$ for smooth moduli,
arXiv:2505.00653v2, 2025.

%% ── Companion papers (Wakil's Program) ───────────────────────────────────────

\bibitem{WakilKhayyam}
K.~Wakil,
Khayyam's triangle: historical restoration and the geometry hidden in the
numbers,
ARC Institute of Knowware, preprint, 2026.

\bibitem{WakilRamanujan}
K.~Wakil,
Ramanujan's dimensional forcing: a universal formula for sieve levels,
ARC Institute of Knowware, preprint, 2026.

\bibitem{WakilEisenstein}
K.~Wakil,
Ternary forcing on the hexagonal lattice: Eisenstein integers, cascade
moduli, and level of distribution~$5/8$,
ARC Institute of Knowware, preprint, 2026.

\bibitem{WakilW3}
K.~Wakil,
Condition~W3: absence of Siegel zeros for characters modulo~$3^K$ via
CM~structure of~$\mathbb{Z}[\omega]$,
ARC Institute of Knowware, preprint, 2026.

\bibitem{WakilShannon}
K.~Wakil,
The Shannon--Wakil effect: constitutional forcing as a universal pattern
in information theory and prime arithmetic,
ARC Institute of Knowware, preprint, 2026.

\end{thebibliography}

\end{document}
